\documentclass[11pt]{amsart}

\usepackage[T1]{fontenc}
\usepackage{lmodern}
\usepackage{microtype}
\usepackage{mathtools,amssymb,amsthm}
\usepackage{upquote}
\usepackage[hidelinks]{hyperref}

\hypersetup{
  pdftitle={An ordered pair of graphs of nonzero twin-width attaining the
            Pettersson-Sylvester tensor-product bound}
}

\newtheorem{theorem}{Theorem}
\newtheorem{lemma}[theorem]{Lemma}
\newtheorem{corollary}[theorem]{Corollary}
\theoremstyle{remark}
\newtheorem{remark}[theorem]{Remark}
\newtheorem*{question}{Question}

\DeclareMathOperator{\tww}{tww}
\newcommand{\sd}{\mathbin{\triangle}}

\title[An ordered pair attaining the tensor-product bound]
{An Ordered Pair of Graphs of Nonzero Twin-Width Attaining
 the Pettersson--Sylvester Tensor-Product Bound}

\date{}

\subjclass[2020]{05C76, 05C75, 68R10}
\keywords{twin-width, tensor product, direct product, circulant graph,
contraction sequence, tightness}

\begin{document}

\begin{abstract}
Pettersson and Sylvester proved that
\[
 \tww(G\times H)\ \le\
 \max\bigl\{\tww(G)\Delta(H)+\Delta(H),\ \tww(H)\bigr\}+\Delta(H)
\]
for the tensor (direct) product of any two graphs, and wrote in their
Conclusion that it would be nice to find two graphs of \emph{nonzero}
twin-width for which this bound is tight. We exhibit such a pair. Let
$G=C_8(1,2)$ be the square of the $8$-cycle and let $H=C_5$; both have
twin-width $2$. Then
\[
 \tww(C_8(1,2)\times C_5)=8=\max\{2\cdot2+2,\ 2\}+2 ,
\]
and more generally $\tww(C_8(1,2)\times C_n)=8$ for every $n\ge5$. The bound is
not symmetric in $(G,H)$, and the claim is for this ordering only. The lower
bound is a three-case count over the $\binom{40}{2}=780$ vertex pairs of the
$40$-vertex product, which is printed here as a graph6 string.
\end{abstract}

\maketitle

\section{The question}

All graphs are finite, simple and undirected. We use the contraction-sequence
definition of twin-width of Bonnet, Kim, Thomass\'e and Watrigant
\cite{TwinwidthI}: a \emph{contraction sequence} of an $n$-vertex graph $X$ is
a sequence of partitions of $V(X)$ starting at the partition into singletons
and ending at the one-part partition, each obtained from the previous one by
merging two parts. A pair of distinct parts $P,Q$ is \emph{black} if every
$p\in P$, $q\in Q$ is an edge, \emph{red} if some such pair is an edge and some
is not, and absent if none is. The \emph{width} of the sequence is the maximum,
over all partitions occurring in it (the initial one included), of the maximum
red degree of a part; $\tww(X)$ is the least width of a contraction sequence of
$X$. Thus $\tww(X)=0$ if and only if $X$ is a cograph.

The \emph{tensor product} (also direct, Kronecker, weak or conjunctive product)
$G\times H$ has vertex set $V(G)\times V(H)$, with $(u,i)(v,j)$ an edge if and
only if $uv\in E(G)$ \emph{and} $ij\in E(H)$. Consequently
\begin{equation}\label{eq:nbhd}
 N_{G\times H}\bigl((u,i)\bigr)=N_G(u)\times N_H(i),
\end{equation}
so $G\times H$ is $\Delta(G)\Delta(H)$-regular whenever $G$ and $H$ are regular.

Pettersson and Sylvester \cite{PS} proved, for arbitrary graphs $G$ and $H$ and
with no hypothesis beyond simplicity,
\begin{equation}\label{eq:PS}
 \tww(G\times H)\ \le\ \max\bigl\{\tww(G)\Delta(H)+\Delta(H),\ \tww(H)\bigr\}
 +\Delta(H).
\end{equation}
The first of the two questions in the unnumbered Conclusion of \cite{PS} reads,
verbatim:
\begin{question}
It would be nice to find an example of two graphs $G,H$ with non-zero
twin-width for which Theorem $[\,$the bound \eqref{eq:PS} above$\,]$ is tight.
\end{question}
\noindent
The bracketed text replaces a cross-reference left unresolved in the source;
nothing else in the sentence is altered. We read ``tight'' as equality in
\eqref{eq:PS} at one ordered instantiation, which is what is exhibited below; a
reading that demands more --- equality for both orderings of the pair, say ---
is not met here. The tight families for \eqref{eq:PS} given in \cite{PS} ---
products with a complete graph, and Hamming graphs --- each have a factor of
twin-width $0$.

\section{The result}

Let $G=C_8(1,2)$ be the circulant graph on $\mathbb Z_8$ in which $u\sim v$
if and only if $u-v\equiv\pm1$ or $\pm2\pmod 8$; equivalently $G$ is the square
of the $8$-cycle. It is $4$-regular with $16$ edges. Let $H=C_5$.

\begin{theorem}\label{thm:main}
$\tww(C_8(1,2))=\tww(C_5)=2$, and
\begin{align*}
 \tww\bigl(C_8(1,2)\times C_5\bigr)
 &=8
 =\max\bigl\{\tww(G)\Delta(H)+\Delta(H),\ \tww(H)\bigr\}+\Delta(H)\\
 &=\max\{2\cdot2+2,\ 2\}+2 .
\end{align*}
\end{theorem}

\begin{corollary}\label{cor:family}
$\tww\bigl(C_8(1,2)\times C_n\bigr)=8=\max\{2\cdot2+2,\ 2\}+2$ for every
$n\ge5$.
\end{corollary}

The active branch of the maximum is $\tww(G)\Delta(H)+\Delta(H)=6>\tww(H)=2$,
so the term attained is the one that carries $\tww(G)\ne0$.

The only tool needed for the lower bounds is \cite[Lemma~3.1]{AHKO}, which is
quoted in \cite{PS} as well.

\begin{lemma}[Ahn, Hendrey, Kim, Oum \cite{AHKO}]\label{lem:lb}
For every graph $X$ with at least two vertices,
\[
 \tww(X)\ \ge\ b_{\min}(X)
 :=\min_{u\ne v}\bigl|\,\bigl(N(u)\sd N(v)\bigr)\setminus\{u,v\}\,\bigr| .
\]
\end{lemma}

We write $b(u,v)=|(N(u)\sd N(v))\setminus\{u,v\}|$ and
$s(u,v)=|N(u)\sd N(v)|$, and $b_{\min}$, $s_{\min}$ for their minima over
unordered pairs of distinct vertices. Lemma~\ref{lem:lb} bounds $\tww$ below by
the \emph{minimum}, so exhibiting one pair with a small $b$ proves nothing: the
whole profile must be computed.

\subsection*{The two factors}

In $C_5$ we have $b(u,v)=2$ for every one of the $10$ pairs (adjacent pairs
have disjoint neighbourhoods, so $s=4$ and two of the four elements are $u,v$
themselves; pairs at distance $2$ have $s=2$ and neither $u$ nor $v$ lies in
the symmetric difference). Hence $b_{\min}(C_5)=2$ and $\tww(C_5)\ge2$. For the
matching upper bound, contract $0$ with $2$, then $3$ with $4$, then the part
$\{0,2\}$ with $1$, then the two remaining parts: the maximum red degree never
exceeds $2$. So $\tww(C_5)=2$.

In $G=C_8(1,2)$, with $N(u)=\{u\pm1,u\pm2\}$, the quantities $s$ and $b$ depend
only on the difference $d=v-u\in\{1,2,3,4\}$:
\[
\begin{array}{c|cccc}
 d & 1 & 2 & 3 & 4\\\hline
 s(u,u+d) & 4 & 6 & 4 & 4\\
 b(u,u+d) & 2 & 4 & 4 & 4
\end{array}
\]
(For $d=1$: $N(0)=\{1,2,6,7\}$ and $N(1)=\{0,2,3,7\}$ have symmetric
difference $\{0,1,3,6\}$, from which $0$ and $1$ are deleted.) Hence
$s_{\min}(G)=4$ and $b_{\min}(G)=2$, so $\tww(G)\ge2$; and the contraction
sequence merging $0$ with $1$, $2$ with $3$, $4$ with $5$, $6$ with $7$, then
$\{0,1\}$ with $\{4,5\}$, then $\{2,3\}$ with $\{6,7\}$, then the two remaining
parts, has width $2$. So $\tww(G)=2$. Both factors have nonzero twin-width, as
the question demands.

\subsection*{The product}

Put $P=G\times H$. By \eqref{eq:nbhd} it has $40$ vertices, is $8$-regular, and
so has $160$ edges. The upper bound $\tww(P)\le8$ is \eqref{eq:PS} itself, so
all the content is in the lower bound, which we now prove by
Lemma~\ref{lem:lb}.

\begin{proof}[Proof that $\tww(P)\ge8$]
Let $(u,i)\ne(v,j)$ and put $m=|N_G(u)\cap N_G(v)|$ and
$m'=|N_H(i)\cap N_H(j)|$. From the profile above, $m\le2$ (indeed
$m=(8-s(u,v))/2$), and in $C_5$ we have $m'=0$ for adjacent $i,j$ and $m'=1$
for $i,j$ at distance $2$. By \eqref{eq:nbhd},
$|N_P((u,i))|=|N_P((v,j))|=8$. Three cases.

\emph{(A) $u=v$, $i\ne j$.} Here
$N_P((u,i))\sd N_P((u,j))=N_G(u)\times(N_H(i)\sd N_H(j))$, and neither $(u,i)$
nor $(u,j)$ lies in this set, since $u\notin N_G(u)$; so nothing is deleted and
$b=4\,|N_H(i)\sd N_H(j)|\ge4\cdot2=8$, with equality exactly when $i,j$ are at
distance $2$ in $C_5$.

\emph{(B) $i=j$, $u\ne v$.} Symmetrically
$b=2\,|N_G(u)\sd N_G(v)|=2\,s(u,v)\ge2\cdot4=8$, with equality exactly when
$uv\in E(G)$.

\emph{(C) $u\ne v$ and $i\ne j$.} Then
$|N_P((u,i))\sd N_P((v,j))|=8+8-2\,|N_G(u)\cap N_G(v)|\cdot|N_H(i)\cap N_H(j)|
=16-2mm'$. Moreover $(u,i)$ lies in the symmetric difference if and only if
$uv\in E(G)$ and $ij\in E(H)$, and likewise for $(v,j)$; so at most $2$ is
deleted, and only when $uv\in E(G)$ and $ij\in E(H)$. If $ij\in E(H)=E(C_5)$
then $m'=0$ because $C_5$ is triangle-free, and $b\ge16-0-2=14$. Otherwise
nothing is deleted and $b=16-2mm'\ge16-2\cdot2\cdot1=12$.

So $b_{\min}(P)=8$. It is attained in case (A), for instance at the pair
$\{(0,0),(0,2)\}$, and in case (B), for instance at $\{(0,0),(1,0)\}$; and
Lemma~\ref{lem:lb} then gives $\tww(P)\ge8$.
\end{proof}

\begin{remark}[the upper bound without \eqref{eq:PS}]\label{rem:noncirc}
Taking $\tww(P)\le8$ from \eqref{eq:PS} is legitimate --- tightness of a bound
\emph{is} the assertion that the bound is attained --- but it leaves the
equality resting on the statement being tested. It need not. The accompanying
program builds an explicit contraction sequence of the $40$-vertex product,
$39$ merges long, and re-certifies its width as $8$ from the original
adjacency; that sequence is printed in the recorded run. So $\tww(P)\le8$ holds
independently of \eqref{eq:PS}, and with $b_{\min}(P)=8$ it brackets
$\tww(P)=8$ without \eqref{eq:PS} entering the argument. The same holds for
each $C_n$, $n=5,\dots,12$.
\end{remark}

This proves Theorem~\ref{thm:main}. For Corollary~\ref{cor:family}, replace
$C_5$ by $C_n$ with $n\ge5$. Then $\Delta(C_n)=2$ and $\tww(C_n)=2$: the lower
bound is $b_{\min}(C_n)=2$ exactly as for $C_5$ (adjacent pairs give
$s=4$ and $b=2$; pairs at distance $2$ give $s=b=2$; pairs at distance $\ge3$
give $s=b=4$), and the upper bound is the round-the-cycle sequence, which merges
$0$ with $1$ and then absorbs $2,3,\dots,n-1$ one at a time and has width $2$.
So the right-hand side of \eqref{eq:PS} is again
$\max\{2\cdot2+2,2\}+2=8$. In the three-case count, (A) gives
$4\,|N_{C_n}(i)\sd N_{C_n}(j)|\ge8$, with equality at distance-$2$ pairs;
(B) gives $2\,s(u,v)\ge8$ unchanged; and in (C) still $m\le2$, while
$m'\le1$ with $m'=0$ whenever $ij\in E(C_n)$ (true for $n\ge5$, as $C_n$ is
then triangle-free), so $b\ge12$. Hence $b_{\min}(C_8(1,2)\times C_n)=8$ and
$\tww(C_8(1,2)\times C_n)=8$. The value $n=4$ is excluded because $C_4$ is a
cograph, so $\tww(C_4)=0$ and $C_4$ is not an admissible factor here.

\medskip
\noindent\textbf{Scope.}
Neither half of the certificate is new: the upper half is \eqref{eq:PS}, which
is the theorem of \cite{PS}, the lower half is \cite[Lemma~3.1]{AHKO}, and the
value $\tww(C_n)=2$ used for the family is also in \cite{AHKO}. We prove no new
lemma; the content here is the pair together with the three-case count. Nothing
beyond equality in \eqref{eq:PS} at the ordered pairs $(C_8(1,2),C_5)$ and
$(C_8(1,2),C_n)$, $n\ge5$, is claimed: with the roles exchanged the bound reads
$\max\{2\cdot4+4,2\}+4=16$ and is not attained by the same product; no
minimality is asserted, neither of the order $40$ nor in any other sense; and we
do not claim that no better general bound of the same shape exists. We looked
for an earlier example of this kind in the nine works recorded as citing
\cite{PS} by OpenAlex and by Semantic Scholar as of 29 August 2026, in the DBLP
title index for twin-width, and in arXiv abstract searches pairing twin-width
with the tensor, direct and Kronecker products, and found none; the nearest
work found was \cite{HIRV}, on graphs whose twin-width equals the bound of
Lemma~\ref{lem:lb}, which does not treat products. That search covered titles,
abstracts and citation records rather than full texts, so we make no priority
claim.

\section{The witness, and reproduction}

Theorem~\ref{thm:main} and Corollary~\ref{cor:family} need no computer: the
product is determined by the two parameters, and the lower bound is the
three-case count above, whose inputs are the two four-entry tables of $s$ and
$b$. For machine reading we also record $P=C_8(1,2)\times C_5$ explicitly.
Label the vertex $(u,i)$ by $k=5u+i$, so that $u=\lfloor k/5\rfloor$ and
$i=k\bmod5$ for $k=0,\dots,39$. Then $P$ has the graph6 string obtained by
concatenating the three lines
\begin{verbatim}
g?@LA_gcASsgSgISQc?Ch@Q_Ad?Ad?HQ??Ch?Ad??IS?
?Sg?AS_??AS_?Ad???Sg??@Q_??QcAO??dL??@Q`O??S
gI??AdC_?@IOHQ??ChhO??SgSg??ISDI??AdCh??AS_
\end{verbatim}
\noindent
(one line of $131$ characters in total, of lengths $44$, $44$ and $43$), whose
$160$ edges are exactly the pairs
$\{5u+i,\,5v+j\}$ with $u-v\equiv\pm1,\pm2\pmod 8$ and $i-j\equiv\pm1\pmod5$.
The only character in the string that is easy to mistranscribe is the single
grave accent (ASCII \texttt{0x60}) on the second line; every other character is
a letter, a digit, \texttt{?}, \texttt{@} or \texttt{\_}.

\medskip
\noindent\textbf{Verification.}
The file \texttt{verify.py} re-derives the computational claims above in exact
integer and set arithmetic using only the Python standard library, and
\texttt{verify.output.txt} records a run of it. It decodes the graph6 string
printed here and checks label-for-label that it is the tensor product built from
the two parameters; recomputes the full $s$- and $b$-profiles of $C_8(1,2)$ and
of $C_5$; re-verifies the two printed contraction sequences at width $2$ in a
from-scratch trigraph simulator, so that $\tww(C_8(1,2))=\tww(C_5)=2$ exactly
with no exhaustive search in the load-bearing path; enumerates all
$\binom{40}{2}=780$ vertex pairs of the product, splits them into the three
cases and reports the per-case minima $8$, $8$, $12$; evaluates both orderings
of \eqref{eq:PS}, obtaining $8$ and $16$; builds and re-certifies the width-$8$
sequence of Remark~\ref{rem:noncirc}; and repeats both bounds for $C_n$,
$n=5,\dots,12$. Its search for a contraction sequence of the $40$-vertex
product is a heuristic, so it proves $\tww(P)\le8$ --- the width of what it
returns is re-certified from the adjacency, which is what makes that sound ---
and it does \emph{not} exclude a width-$7$ sequence; only Lemma~\ref{lem:lb}
does that. Nothing in the proof above needs the program: every input it consumes
is printed here.

\begin{thebibliography}{9}

\bibitem{AHKO}
J.~Ahn, K.~Hendrey, D.~Kim, and S.~Oum,
\emph{Bounds for the twin-width of graphs},
SIAM J. Discrete Math. \textbf{36} (2022), no.~3, 2352--2366.
\href{https://doi.org/10.1137/21M1452834}{doi:10.1137/21M1452834};
\href{https://arxiv.org/abs/2110.03957}{arXiv:2110.03957}.

\bibitem{TwinwidthI}
\'E.~Bonnet, E.~J. Kim, S.~Thomass\'e, and R.~Watrigant,
\emph{Twin-width I: tractable FO model checking},
J. ACM \textbf{69} (2022), no.~1, Art.~3, 46 pp.
\href{https://doi.org/10.1145/3486655}{doi:10.1145/3486655};
\href{https://arxiv.org/abs/2004.14789}{arXiv:2004.14789}.

\bibitem{HIRV}
F.~Heinrich, F.~Ihringer, M.~Rassmann, and T.~Volk,
\emph{On the twin-width of near-regular graphs},
Discrete Appl. Math. \textbf{379} (2026), 177--193.
\href{https://doi.org/10.1016/j.dam.2025.07.044}{doi:10.1016/j.dam.2025.07.044};
\href{https://arxiv.org/abs/2504.02342}{arXiv:2504.02342}.

\bibitem{PS}
W.~Pettersson and J.~Sylvester,
\emph{Bounds on the twin-width of product graphs},
Discrete Math. Theor. Comput. Sci. \textbf{25} (2023), no.~1, Paper No.~18.
\href{https://doi.org/10.46298/dmtcs.10091}{doi:10.46298/dmtcs.10091};
\href{https://arxiv.org/abs/2202.11556}{arXiv:2202.11556v4}.

\end{thebibliography}

\end{document}
